\section{IA, energía y escala civilizatoria}

La sección anterior aterrizó la comercialización en el costo de la planta eléctrica y en la ruta de combustible; esta sección regresa a la pregunta sobre la IA y la civilización con la que abrió la entrevista. ¿Por qué un emprendedor de la fusión nuclear habla de AGI, centros de datos y capacidad de cómputo? Porque la energía no es una aplicación posterior, sino la base de todo cómputo, toda manufactura y toda organización social; cuando la IA sigue creciendo de forma exponencial, el suministro de energía deja de ser una variable de fondo y se convierte en una restricción de primer plano.

Esta sección sitúa la fusión nuclear en la era de la IA. Yang Zhao (杨钊) sostiene que todo lo que crece de forma exponencial termina por encontrar un cuello de botella; a corto plazo, los cuellos de botella de la IA pueden ser la capacidad de cómputo, los chips, los datos y la energía, pero a largo plazo el suministro de energía será el verdadero gran cuello de botella. Mientras se pueda ofrecer energía más barata, la demanda casi con certeza la consumirá, del mismo modo que, en cuanto mejora el rendimiento de cómputo, siempre aparecen nuevas aplicaciones que lo agotan.

A la inversa, la IA también puede ayudar a la fusión nuclear. En la entrevista se mencionan tres funciones: primero, la IA puede usarse para el control en tiempo real y sustituir modelos tradicionales complejos pero demasiado lentos; segundo, la IA puede mejorar el diagnóstico y reemplazar con algoritmos parte del hardware costoso; tercero, la IA puede entrenar, con datos experimentales reales, modelos de simulación a nivel de dispositivo y acortar el ciclo de iteración de los experimentos. Es decir, la relación entre la IA y la fusión no es unidireccional: la IA necesita energía y la fusión también necesita a la IA para reducir costos y ganar eficiencia.

\begin{figure}[H]
\centering
\includegraphics[width=0.94\textwidth]{energy-ai-loop.png}
\caption{La IA y la energía se amplifican mutuamente: el crecimiento de la capacidad de cómputo eleva la demanda de electricidad y la energía de bajo costo, a su vez, eleva el techo de la civilización. Diagrama conceptual propio, elaborado a partir de la conversación entre 01:44:42 y 01:50:00.}
\end{figure}

\begin{knowledgebox}{Lectura de la figura: la energía no es un trasfondo externo de la IA}
La figura no se limita a conectar Fusion con Data Center, sino que muestra una retroalimentación positiva: la IA consume más energía e impulsa la inversión en nuevas fuentes de energía; las nuevas fuentes reducen los costos y, a su vez, hacen posible más cómputo, más manufactura y más experimentos. Si la fusión nuclear tiene éxito, será uno de los amplificadores de base de la era de la IA.
\end{knowledgebox}

\begin{knowledgebox}{Tres tipos de ayuda de la IA a la fusión}
\begin{tabularx}{\textwidth}{p{0.20\textwidth}p{0.34\textwidth}X}
\toprule
Escenario & Qué hace & Valor \\
\midrule
Control en tiempo real & Usa IA para aproximar procesos físicos complejos y entregar con rapidez estrategias de control & Reduce la latencia de control y mejora la estabilidad del plasma. \\
Diagnóstico mejorado & Usa algoritmos para aumentar la precisión de la observación o sustituir equipos de diagnóstico costosos & Reduce el costo de hardware y mejora la calidad de la estimación del estado. \\
Simulación experimental & Entrena, con datos experimentales reales, un surrogate model a nivel de dispositivo & Reduce el número de pruebas de ensayo y error y acelera el ajuste de parámetros y la optimización de la operación. \\
\bottomrule
\end{tabularx}
\end{knowledgebox}

\subsection{Seguridad, regulación y elección del combustible}

Esta subsección completa el tema regulatorio que toda planta eléctrica comercial debe enfrentar. La subsección anterior dijo que la IA puede reducir costos y ganar eficiencia, pero un producto energético no es un modelo de laboratorio: el combustible, los residuos y la regulación se convierten directamente en costos. Yang Zhao distingue en particular entre deuterio-tritio y deuterio-deuterio porque su dificultad física y su dificultad comercial van en sentidos opuestos: el deuterio-tritio reacciona con más facilidad, mientras que el deuterio-deuterio se acerca más a la visión de largo plazo de baja regulación y poca presión sobre el combustible.

Una planta de fusión no solo necesita una Q alta; también debe ser regulable, mantenible y aceptada por la sociedad. En la entrevista se mencionó en particular la diferencia entre la ruta deuterio-tritio y la ruta deuterio-deuterio. El tritio puede usarse en sistemas relacionados con armas, está sujeto a una regulación estricta, es costoso y difícilmente existe en la naturaleza de forma estable y en grandes cantidades, por lo que hay que generarlo por reproducción mientras se produce electricidad; esto aumenta la complejidad de ingeniería y el costo regulatorio. Energy Singularity (能量奇点) aspira a largo plazo a la fusión deuterio-deuterio, porque las reservas de deuterio en el agua de mar son enormes y ofrecen ventajas en disponibilidad de materia prima, seguridad y manejo de residuos, aunque esta ruta requiere parámetros más altos.

\begin{warningbox}{La ruta de combustible no solo afecta la dificultad física}
Con deuterio-tritio es más fácil alcanzar las condiciones de reacción, pero la reproducción del tritio, la regulación y los sistemas de seguridad aumentan la complejidad de la planta; el deuterio-deuterio es más difícil, pero si se logra, la presión sobre la materia prima y la regulatoria podría disminuir de forma notable. La ruta de comercialización debe comparar a la vez la dificultad física y el costo de la planta completa.
\end{warningbox}

\begin{knowledgebox}{Comparación de rutas de combustible}
\begin{tabularx}{\textwidth}{p{0.18\textwidth}p{0.34\textwidth}X}
\toprule
Ruta & Ventajas & Dificultades \\
\midrule
Deuterio-tritio & Alcanza con más facilidad las condiciones de reacción; es la ruta de validación predominante a corto plazo & El tritio está regulado, requiere reproducción y conlleva una fuerte presión en seguridad del sistema y en costos. \\
Deuterio-deuterio & El deuterio proviene del agua de mar; ventajas claras de materia prima y de regulación a largo plazo & La reacción es más difícil; requiere un producto triple más alto y un dispositivo con mejor desempeño. \\
Planta de demostración & Puede demostrar primero que Q, la extracción de calor y la conexión a la red son viables & Todavía no equivale a la ruta final de menor costo y menor regulación. \\
\bottomrule
\end{tabularx}
\end{knowledgebox}

\subsection{Qué ocurre después de la energía ilimitada}

Esta subsección lleva la discusión de la economía de la planta eléctrica a la escala civilizatoria. Cada paso anterior es muy ingenieril: imanes, Q, costos, regulación, financiamiento; pero si Yang Zhao está dispuesto a tratar este asunto como algo «para ver en vida», es porque, una vez que cruce el umbral de la comercialización, ya no se limitará a sustituir una fuente de electricidad, sino que ampliará el presupuesto energético que la humanidad puede costear.

Cuando el costo de la energía baje de forma drástica, muchas cosas que hoy «no son rentables» se volverán viables: la desalinización de agua de mar, el cómputo a gran escala, la fabricación de materiales, la ingeniería en entornos extremos y la propulsión interplanetaria, entre otras, tendrán nuevos límites. Yang Zhao considera la fusión un cambio a escala civilizatoria: no solo sustituye a la generación termoeléctrica, sino que podría hacer que la manera en que la humanidad usa la energía salte, en conjunto, un orden de magnitud.

\begin{importantbox}{El significado civilizatorio del suministro de energía}
La energía no es el insumo de una sola industria, sino la restricción de base del cómputo, la manufactura, el transporte, las ciudades, la agricultura y las actividades espaciales. Si la fusión logra un suministro de energía estable y de bajo costo, lo que cambia es el conjunto de cosas que la civilización puede hacer.
\end{importantbox}

\begin{warningbox}{No hay que entender «energía ilimitada» como ausencia de restricciones}
Aun si la fusión tiene éxito, siguen existiendo la construcción de equipos, los materiales, la operación y el mantenimiento, la regulación, la transmisión eléctrica, el despacho de la red y los ciclos de capital. Es más preciso decir que la fusión podría reducir de forma notable el costo marginal de la energía y las restricciones de recursos, pero no hará desaparecer automáticamente todas las restricciones de ingeniería.
\end{warningbox}

\subsection{Riesgo de fracaso y el modelo de escalar la montaña}

La subsección anterior trató el espacio civilizatorio que se abre tras el éxito energético; esta subsección vuelve al problema real que el emprendedor enfrenta cada día: cómo escalar esta montaña y en qué puntos podría caer. La visión grandiosa de la fusión nuclear hace fácil pasar por alto la dureza del nivel de ejecución, pero el juicio final de Yang Zhao es claro: con suficiente tiempo y dinero, la ruta tiene solución desde el punto de vista de la ingeniería; el verdadero riesgo es no poder poner en marcha el dispositivo, no poder entregarlo o que el equipo no logre seguir elevando sus propias capacidades.

Al cierre de la entrevista, Yang Zhao compara emprender en fusión con escalar montañas de forma continua, y la altura de las montañas crece de manera exponencial. Las causas de fracaso pueden condensarse en tres tipos: falta dinero y el siguiente dispositivo no puede ponerse en marcha; falta gente y el juicio y la ejecución de ingeniería en los puntos decisivos no están a la altura; el trabajo no se concreta y el dispositivo no alcanza sus objetivos de diseño, de modo que el 380 posterior no podrá venderse. Es un juicio muy sereno: el PMF de la fusión no es ambiguo, la demanda es real y el producto es claro; el verdadero riesgo está en si se logra construir el producto.

\begin{knowledgebox}{Desglose de riesgos: por qué no es un emprendimiento común}
\begin{tabularx}{\textwidth}{p{0.18\textwidth}p{0.34\textwidth}X}
\toprule
Riesgo & Manifestación concreta & Consecuencia \\
\midrule
Falta dinero & El 170 o un subsistema esencial no puede ponerse en marcha & La ruta tecnológica se queda detenida en la montaña anterior. \\
Falta gente & No se toman decisiones de ingeniería correctas cuando la información es insuficiente & El avance se retrasa y los errores son difíciles de localizar. \\
El trabajo no se concreta & Los parámetros entregados no alcanzan los objetivos de diseño & Falla la lógica de venta del 380 y de la planta eléctrica comercial. \\
\bottomrule
\end{tabularx}
\end{knowledgebox}

\subsection{Resumen de la sección}

El crecimiento de la IA a largo plazo necesita energía, y la investigación y el desarrollo de la fusión nuclear también pueden acelerarse con IA. La comercialización de la fusión no solo significa sumar una nueva forma de generar electricidad, sino que podría cambiar los límites del cómputo, la industria y la expansión de la civilización.
